\chapter{Runtime Systems}

\section{Rendering}

The renderer is implemented in \filepath{src/demi/runtime/Renderer2D.cpp}. It uses SDL3 when available. Current renderer capabilities include:

\begin{itemize}
  \item clearing with the active camera color;
  \item drawing PPM-backed textures when available;
  \item drawing fallback rectangles for entities without textures;
  \item drawing collider outlines for debug visibility;
  \item drawing debug lines;
  \item drawing HUD rectangles, buttons, and pixel text;
  \item scaling HUD coordinates from each HUD file canvas to the active window size.
\end{itemize}

The HUD text renderer is a built-in pixel font. It currently supports the glyphs needed by the example menu and HUD.

\section{Physics}

The 2D physics system is implemented in \filepath{src/demi/runtime/Physics2D.cpp}. It supports:

\begin{itemize}
  \item dynamic and static rigidbodies;
  \item gravity scale per dynamic body;
  \item axis-aligned box colliders;
  \item static collider resolution;
  \item bounciness on vertical contacts;
  \item overlap box queries exposed to Lua;
  \item regression coverage for edge-contact teleport behavior.
\end{itemize}

The resolver uses previous collider bounds for side decisions, which avoids incorrectly snapping a player across a platform when touching a top-right or top-left edge.

\section{Audio}

The audio system is implemented in \filepath{src/demi/runtime/AudioSystem.cpp}. It supports:

\begin{itemize}
  \item asset-backed audio loading;
  \item MP3 decoding with libmpg123 when available;
  \item playing and stopping audio handles from Lua;
  \item master volume control;
  \item mixing currently playing sounds into the SDL audio callback.
\end{itemize}

\section{Window and Runtime Control}

The runtime supports Lua-requested window mode changes:

\begin{itemize}
  \item \code{windowed};
  \item \code{borderless};
  \item \code{fullscreen}.
\end{itemize}

Window mode changes are requested from Lua and applied by the runtime loop. Lua can also pause or resume physics through \code{Runtime.set\_physics\_enabled} and request shutdown through \code{Runtime.quit}.

\section{Scene Switching}

Runtime scene switching is implemented in \filepath{src/demi/runtime/SceneLoader.cpp} and \filepath{src/demi/runtime/LuaScriptHost.cpp}. The free function \code{loadScene} resolves a scene id against the project scene list, validates the scene and optional HUD, and returns a parsed \code{World}.

Lua calls \code{Scene.load} from any callback. The host records the request as pending. The application loop in \filepath{src/demi/runtime/RuntimeApp.cpp} checks for a pending scene after the Lua update step and, if present, rebuilds the active world, destroys and re-loads scripts, and re-runs \code{on_create}/\code{on_start}. Switches are applied between frames so they never re-enter Lua mid-callback.
